\documentclass[10pt,a4paper]{amsart}
\usepackage[margin=19mm]{geometry}
\usepackage{amsmath,amssymb,mathtools}
\usepackage[scr=rsfs,cal=euler]{mathalfa}
\usepackage{xfrac}
\usepackage[unicode,hidelinks,bookmarks=false]{hyperref}
\newcommand{\ie}{i.e.\ }
\newcommand{\cf}{cf.\ }
\newcommand{\resp}{respectively\ }
\newcommand{\Subsection}{\subsection}
\providecommand{\nobreakdash}{\nobreak\hbox{-}}
\setlength{\emergencystretch}{2em}
\multlinegap0pt
\usepackage{amssymb}
\usepackage[shortlabels]{enumitem}
\newtheorem{theorem}{Theorem}
\title{Families of Unit Equations and Exponential Diophantine Problems via Integral Points}
\author{Julie Tzu-Yueh Wang, Zheng Xiao}
\date{Source: arXiv:2604.26497v1 (CC BY 4.0)}
\begin{document}
\maketitle

\noindent\emph{Source and licence.} Families of Unit Equations and Exponential Diophantine Problems via Integral Points. Version arXiv:2604.26497v1 (CC BY 4.0), distributed by arXiv under the Creative Commons Attribution 4.0 International licence, http://creativecommons.org/licenses/by/4.0/, as recorded in the arXiv licence metadata for this exact version and shown on the abstract page https://arxiv.org/abs/2604.26497v1. Full licence notice: \texttt{licenses/arxiv-2604.26497-CC-BY-4.0.txt}.\par\smallskip

\noindent\textit{Editorial scope.} Counted unit: the article's theorem UnitCZ (source0.tex lines 670--676). The article's complete original proof of it is delivered with it (source0.tex lines 678--681, one \texttt{proof} environment of 4 lines). No mathematics is written, repaired or re-derived: the statement and the proof are the article's own lines, in the article's own notation, and no retained line is substituted or re-flowed.  Retained uncounted as the source-local context the proof needs: lines 373--381.  Nothing was omitted from any retained span, so the package carries no omission record.  No retained line contains a citation, so the wrapper carries no bibliography and no bibliography entry.  No retained line contains a figure environment, an image-inclusion command or drawing code, so no image is delivered and the package declares no asset dependency.  Editorial formatting only, in addition: an AMS \texttt{amsart} preamble that restates the article's own declarations and macros used by the retained lines, namely the environments \texttt{theorem} on separate counters, together with the standard packages those lines need.  The retained environments print in this document as Theorem 1 in order of appearance, whereas the article prints the same items with the numbering of its own full document.  The complete native source of the article as distributed in the arXiv e-print for this exact version is bundled unchanged as \texttt{sources/199399/source0.tex}.
\begin{theorem}\label{unit2}
Let $f_i\in k[t]$, $1\le i\le n+1$, be non-constant polynomials of the same degree $d$.
Let  $V$  be the hypersurface  in $\mathbb A^{n+1}$ defined by the zero locus of  
   \begin{equation}\label{funiteq}
        f_1(x_0)x_1 + \cdots + f_n(x_0)x_n=f_{n+1}(x_0).
    \end{equation} 
  Then there exists a proper Zariski closed subset $Z\subset V$    such that all but finitely many solutions $(t,u_1,\hdots,u_{n})\in \mathcal O_S\times  (\mathcal O_S^*)^n$ of the equation  \eqref{funiteq}
are contained in $Z$.  
 \end{theorem}

\begin{theorem}\label{UnitCZ}
Let $f=\sum_{i=1}^n x_i g_i$, where $g_i \in k[x_0,\ldots,x_n]$, $n\ge 3$, are homogeneous polynomials of the same degree such that $g_i(1,0,\ldots,0)\ne 0$, and  let $X$  be the hypersurface  in $\mathbb P^n$ defined by the zero locus of  $f=\sum_{i=1}^nx_ig_i$.   Then there exists a Zariski-closed subset $Z \subset X$ such that all but finitely many solutions of
\[
f(y,u_1,\ldots,u_n)=0, \quad (y,u_1,\ldots,u_n)\in \mathcal O_S\times \mathcal O_S^*\times\cdots\times \mathcal O_S^*
\]
are contained in $Z$.
\end{theorem}

\begin{proof}[Theorem \ref{UnitCZ} implies  Theorem  \ref{unit2}]
Let $g_i=x_{n+1}^d f_i(\frac{x_0}{x_{n+1}})$ for $1\le i\le n$ and let $g_{n+1}= -x_{n+1}^d f_{n+1}(\frac{x_0}{x_{n+1}}) $. These are homogeneous polynomials of degree $d$.  Furthermore, because each $f_i$ has degree exactly $d$, which implies
 $g_i(1,0,\hdots,0)\ne 0$ for all  $1\le i\le n+1$.   Let $X$  be the hypersurface  in $\mathbb P^{n+1}$ defined by the zero locus of  $F=\sum_{i=1}^{n+}x_ig_i $.  Theorem  \ref{unit2} then follows directly from Theorem \ref{UnitCZ} by considering the affine patch where $x_{n+1}=1$.
\end{proof}

\end{document}
